% PNP RESEARCH REPORT::COMPATIBILITY-CORE-DEGREE-TRANSFER::2026-11-11T20:35:00-08:00
\documentclass[11pt]{article}
\usepackage[T1]{fontenc}
\usepackage[utf8]{inputenc}
\usepackage[margin=1in]{geometry}
\usepackage{amsmath,amssymb,amsthm,booktabs,longtable,array,xurl}
\usepackage[hidelinks]{hyperref}
\setlength{\emergencystretch}{3em}
\newtheorem{theorem}{Theorem}
\newtheorem{proposition}{Proposition}
\newtheorem{corollary}{Corollary}
\newcommand{\F}{\mathbb F_2}
\newcommand{\poly}{\operatorname{poly}}
\newcommand{\SAT}{\operatorname{SAT}}
\newcommand{\degM}{\operatorname{deg}_{\mathcal M}}
\title{Compatibility Constraints Revisited:\\
An Exact Degree Transfer and a Stronger Selector Algorithm}
\author{P vs NP Research Project\\Codex research with independent temporary-agent review}
\date{11 November 2026, 20:35 PST}
\begin{document}
\maketitle

\section*{Status and new attempt}
\textbf{P versus NP remains open.}
The user requested another attempt after the adaptive collision and
closure barriers \cite{barriers2026}. This report deliberately removes
the all-prefix-extension premise. It constructs and checks formulas
with locally feasible prefixes that have no global continuation.

Three exact results emerge. First, the recovered selector encoding has
a minimum restricted algebraic degree equal to the source clause count
on satisfiable inputs and zero on unsatisfiable inputs.
Second, this special encoding admits a global $2^m\poly(N)$ controller,
stronger than the generic $N^{O(m)}$ closure-seed bound.
Third, adaptive collisions from efficiently invertible permutations have
explicit global coordinates, even after a final compatibility restriction.
These are \textbf{PROVED/EXACT} within the stated models; finite checks
are \textbf{COMPUTATIONAL}. The all-program work amplification needed
for a separation remains \textbf{OPEN}.

\section*{Quantifiers, input measure, and SAME\_REALIZATION}
Inputs are explicit Boolean formulas or circuits with dense occurring
variable names, or with a charged preprocessing step making them so.
$N$ denotes total encoded input length plus two, including the source,
compatibility conditions, construction metadata, and any explicitly
supplied stage circuits. The source clause count is $m$.
For a fixed $q\ge1$, a degree-$q$ quasipolynomial bound means
$2^{O((\log N)^q)}$; this is an upper bound unless stated otherwise.

The previous minimum-work benchmark is
\[
 H(F)=\min_{|p|\le b_0+\lfloor\log(|F|+2)\rfloor}
       \{t_p(F)+q_0(|F|+2)+1\},
 \tag{1}
\]
where successful outputs are independently verified witnesses and
$q_0$ is a fixed polynomial verification allowance.
One fixed correct SAT decider of worst-case time $T_A$ gives
\[
 \overline H(N)\le C_A(N+T_A(KN)+2)^D
 \tag{2}
\]
for fixed $K,D$. The preceding transfer report proves the converse
polynomial and fixed-quasipolynomial-degree equivalences
\cite{transfer2026}.
The all-program amplification target is
\[
 s_{j+1}\le Cs_j,\qquad
 H(F_{j+1})\ge H(F_j)s_j^c/B,\quad s_j\longrightarrow\infty,
 \tag{3}
\]
with fixed $C>1,c>0,B\ge1$. It would force
$H(F_j)\ge2^{\Omega((\log s_j)^2)}$.
None of the structural results below establishes (3).
The degree resource and the witness-work resource are not identified.

\section{An exact compatibility encoding recovered from the corpus}
Let
\[
 F(x)=\bigwedge_{j=1}^m C_j(x),\qquad
 C_j(x)=\bigvee_{t=1}^{k_j}L_{jt}(x),
 \tag{4}
\]
with $m\ge1$ and every clause nonempty.
Introduce one selector $s_{jt}$ per literal occurrence.
The Horn 2-CNF is
\[
 B_F(x,s)=
 \bigwedge_j
 \left[
 \bigwedge_{t<u}(\neg s_{jt}\vee\neg s_{ju})
 \ \wedge\
 \bigwedge_t(\neg s_{jt}\vee L_{jt}(x))
 \right].
 \tag{5}
\]
Add the disjoint affine rows
\[
 \bigoplus_{t=1}^{k_j}s_{jt}=1,\qquad j=1,\ldots,m.
 \tag{6}
\]
Pair exclusions enforce at most one selected literal, and odd parity
enforces at least one. Thus exactly one true literal is selected in
every clause:
\[
 \boxed{
 F(x)\ \Longleftrightarrow\
 \exists s\,[B_F(x,s)\wedge A_Fs=\mathbf1].
 }
 \tag{7}
\]
Both witness transformations are explicit. A source model selects one
true literal per clause; a target model projects to its $x$ assignment.

This is the actual construction in AFFINE-TABLE-COMPOSITION2 and in its
\texttt{selector\_encoding} implementation \cite{atc2026,atccode2026}.
For exact 3-CNF it uses $n+3m$ variables, $6m$ Horn binary clauses,
and $m$ disjoint nonzero affine rows of rank exactly $m$.
Here $n$ counts occurring original variables.
For general widths, the Horn clause count is
$\sum_j[\binom{k_j}{2}+k_j]$, still polynomial in explicit source length.
An empty source clause is handled separately as immediate inconsistency.

\section{Exact degree dichotomy: the algebraic property encodes SAT}
Let $\mathcal M_F$ be the set of all Horn models of $B_F$,
identified with Boolean vectors $(x,s)$. This family is nonempty:
all selectors zero satisfies the Horn part for every $x$.
It is also closed under intersection.

The affine acceptance function on $\mathcal M_F$ is
\[
 a_F(x,s)=
 \prod_{j=1}^m\left(\sum_{t=1}^{k_j}s_{jt}\right)
 \quad\text{over }\F.
 \tag{8}
\]
Its value is one precisely when all rows (6) hold.
Define
\[
 D(F)=\min\{\deg p:
 p\text{ is a multilinear }\F\text{-polynomial agreeing with }
 a_F\text{ on }\mathcal M_F\}.
 \tag{9}
\]
Use zero as a nonnegative degree bound for the zero function.

\begin{theorem}[Exact restricted-degree transfer]\label{thm:degree}
For every source formula in (4),
\[
 \boxed{
 D(F)=
 \begin{cases}
  m,& F\text{ is satisfiable},\\
  0,& F\text{ is unsatisfiable}.
 \end{cases}}
 \tag{10}
\]
\end{theorem}

\begin{proof}
If $F$ is unsatisfiable, (7) says no Horn model satisfies every affine
row. Thus $a_F$ is identically zero on $\mathcal M_F$, proving
$D(F)=0$.

Suppose $F$ has a satisfying assignment $a$.
For each clause choose one position $t_j$ whose literal is true under $a$.
Fix $x=a$ and all selectors other than $s_{j,t_j}$ to zero.
Leave the $m$ chosen selectors free, calling them $u_1,\ldots,u_m$.
Every point of this entire Boolean cube satisfies $B_F$:
there is at most one selector per clause, and every possibly selected
literal is true under $a$.
On this cube, (8) becomes
\[
 a_F=\prod_{j=1}^m u_j=\operatorname{AND}_m(u).
\]
Any representing polynomial in (9), after this substitution, must agree
with $\operatorname{AND}_m$ on the full Boolean cube.
The unique multilinear representation of $\operatorname{AND}_m$
has degree $m$, and substitution of constants cannot increase degree.
Therefore $D(F)\ge m$. Formula (8) already gives a degree-$m$
representation, so equality follows.
\end{proof}

The proof does not require agreement outside $\mathcal M_F$:
the lower-bound cube lies wholly inside that model family.

\begin{corollary}[Computational consequence]\label{cor:degree}
On canonical selector-image inputs with their source metadata,
deciding $D(F)<m$ is coNP-complete. Even deciding $D(F)=0$
has the same property. Computing $D(F)$ in polynomial time is
equivalent to polynomial-time SAT.
The equivalence also preserves every fixed quasipolynomial
logarithmic degree.
\end{corollary}

\begin{proof}
By (10), both decision predicates are exactly source UNSAT.
The compiler is polynomial time; validity of the canonical construction
can be checked from the supplied source. On malformed representations
use a fixed rejection convention.
On valid images, the complement has an NP certificate consisting of
a source model; malformed inputs are recognized directly in polynomial time.
A SAT algorithm computes $D(F)$ by outputting $m$ or zero.
Conversely a degree algorithm decides SAT by testing for output $m$.
The reductions have polynomial encoded size, and deterministic
quasipolynomial log-degree is preserved under polynomial substitutions
and complementation.
\end{proof}

This is a transfer of \emph{computational difficulty of determining a
degree property}. It is not a transfer from a large numerical degree
to large computation time.
For example,
\[
 F_m=\bigwedge_{j=1}^m(a\vee u_j\vee v_j)
 \tag{11}
\]
is a connected family with the immediate model $a=1$.
Its selector representation still has $D(F_m)=m$ and rank $m$.
Its witnesses are uniformly easy.

There is a further exact obstruction for the closure-enumeration
procedure. No positive Horn implication in (5) has a selector as its
head. Thus Horn propagation never creates a selector not already
in the seed. Any accepted closure must contain one selector from each
of the $m$ clauses and consequently needs at least $m$ selector seeds.
This proves a limitation of that particular seed-enumeration method,
while (11) has a direct polynomial witness algorithm.

Under $\mathsf P=\mathsf{NP}$, the degree problem in
Corollary \ref{cor:degree} is in P. There is no contradiction with
$D(F)=m$: $\operatorname{AND}_m$ itself has degree $m$ and a linear
evaluation algorithm. A lower bound on \emph{computing} the degree
property would be the missing ingredient, and would already separate
P from NP on this family.

\section{A stronger global controller: exponential in clauses, not seeds}
For original variable $x_i$, let
\[
 P_i=\{j:x_i\in C_j\},\qquad
 M_i=\{j:\neg x_i\in C_j\}.
\]
A dynamic-programming state records the set of clauses already made
true. Put
\[
 \mathcal D_0=\{\varnothing\},
\]
\[
 \mathcal D_i=
 \{S\cup P_i:S\in\mathcal D_{i-1}\}
 \ \cup\
 \{S\cup M_i:S\in\mathcal D_{i-1}\}.
 \tag{12}
\]

\begin{theorem}[Clause-coverage controller]\label{thm:dp}
For an explicitly represented $m$-clause CNF, (12) decides
satisfiability and finds a model in $2^m\poly(N)$ time.
It therefore solves the canonical selector representation (5)--(6)
within the same bound, including construction of a full selector witness.
\end{theorem}

\begin{proof}
Inductively, $\mathcal D_i$ is exactly the family of clause sets made
true by assignments to the first $i$ variables. The two successors
in (12) set $x_i$ to one or zero. Hence the full clause set $[m]$
belongs to $\mathcal D_n$ exactly when a satisfying assignment exists.
There are at most $2^m$ states and two transitions per state per
variable. Bitmask operations and input preparation cost a polynomial
in explicit length. Predecessor records reconstruct a source model,
which is lifted by (7).
Unused variables can first be removed with their identifiers recorded;
the theoretical input convention is dense occurring variables.
\end{proof}

In this special representation $d=\operatorname{rank}A_F=m$, so the
upper bound is
\[
 2^d\poly(N),
 \tag{13}
\]
rather than the generic Horn/affine closure bound $N^{O(d)}$.
The generic tractability results are established in the literature
\cite{brakensiek2019}; (12) is an elementary direct algorithm for the
particular recovered image, with a proof included here.
No priority claim is made for clause-parameter dynamic programming.

For fixed width three, the closure universe has
$U=n+3m\le6m$, since only occurring original variables are counted.
The previous refined binomial bound already yields
$\sum_{j\le m}\binom Uj\le2^U=2^{O(m)}$.
Thus (13) improves the explicit base and controller, not the qualitative
clause-exponential frontier in this constant-width case.
Its distinction from the general seed bound is stronger for unbounded
clause width, where $U$ can be much larger than $m$.

Thus $m=O(\log N)$ is polynomial for this compatibility core.
For an extension chain with polynomially many \emph{actually satisfiable}
selector residuals and uniformly charged representation acquisition,
lifting, and replay, let $m_*$ be the maximum residual clause count.
Then
\[
 H(F_J)\le \poly(N)2^{O(m_*)}.
 \tag{14}
\]
The fixed-program and canonical-reconstruction requirements from the
previous report remain in force.
A squared-logarithm $H$ lower bound on such a chain would require
\[
 m_*=\Omega((\log N)^2).
 \tag{15}
\]
If a previously chosen witness leaves an unsatisfiable residual,
(14) cannot simply be applied to that residual as a successful
extension. One must solve the represented global compatibility problem
or give a charged repair procedure. Equation (12) does solve a fully
supplied global CNF, with its \emph{total} clause count charged.

\section{A concrete locally feasible prefix with no continuation}
Consider
\[
\begin{aligned}
 F={}&(x\vee y\vee z)\\
 &\wedge(\neg x\vee a\vee b)
 \wedge(\neg x\vee a\vee\neg b)\\
 &\wedge(\neg x\vee\neg a\vee b)
 \wedge(\neg x\vee\neg a\vee\neg b).
\end{aligned}
 \tag{16}
\]
The last four clauses force $x=0$. The formula has 12 models, because
$y\vee z$ must hold while $a,b$ are free.
Selecting $x$ in the first clause is locally feasible under the Horn
constraints and first parity row. It has no full continuation.

The new controller finds, for example,
$(x,y,z,a,b)=(0,0,1,0,0)$.
In the recorded run it uses 26 attempted transitions and at most six
states in a layer.
The selector representation has 20 variables, 30 Horn binary clauses,
and five affine rows of rank five.
Its acceptance degree is exactly five by Theorem \ref{thm:degree}.
Therefore dead prefixes, nonconstant degree, and full row rank can
coexist with a small global computation.

\section{Global coordinates for compatible adaptive collision chains}
At stage $i$, let $P_{a_i}$ be a uniformly efficiently invertible
permutation on $r_i+1$ bits. Its key $a_i$ is a deterministic,
efficiently computed function of earlier ordered collision pairs.
Stage validity requires $u_i\ne v_i$ and equal first $r_i$ output bits.

\begin{theorem}[Collision coordinate normal form]\label{thm:coords}
The full chain of ordered collision pairs is in polynomial-time
bijection with coordinates
$(z_1,b_1,\ldots,z_k,b_k)$, where $|z_i|=r_i$ and $b_i$ is one bit:
\[
 u_i=P_{a_i}^{-1}(z_i b_i),\qquad
 v_i=P_{a_i}^{-1}(z_i(1-b_i)).
 \tag{17}
\]
All sizes, key generation, forward/inverse evaluation and stage counts
must be uniformly polynomially charged.
\end{theorem}

\begin{proof}
Given a collision, injectivity says the full outputs differ. Since
their first $r_i$ bits agree, their last bits are complementary.
The common prefix uniquely defines $z_i$, and the last bit of
$P_{a_i}(u_i)$ uniquely defines $b_i$. This proves the single-stage
bijection. In a chain, previous pairs determine the next key.
Sequential application gives the global bijection and its efficient
inverse. Gate values are completed afterward.
\end{proof}

Write $E(z,b)$ for this encoding.
An arbitrary explicit compatibility circuit $K$ on the whole chain
becomes the exact remaining constraint
\[
 K(E(z,b))=1.
 \tag{18}
\]
This is a polynomial-size ordinary circuit under the hypotheses.
It is a normal form, not a polynomial algorithm for arbitrary $K$.

When compatibility is a source CNF $F(z)$ on the common output labels,
the search problems are witness-preserving equivalent.
A source model lifts by choosing every $b_i=0$ in (17).
A compatible chain projects by reading its common output labels.
Each source model has exactly $2^k$ ordered primary-chain lifts
when orientations are unrestricted.

In particular, some forward-chosen prefixes can be dead even though
every satisfying \emph{source assignment} has an easy lift.
A global source-first algorithm avoids the bad choices.
Iterating this wrapper from one fixed witnessed seed, with polynomial
total reconstruction and generation cost in final size, still gives
a fixed short replay program and polynomial $H$.
No pointwise equality of optimal short-program costs is asserted.

\subsection*{A precise inversion alternative}
Remove efficient inverse evaluation. For a fixed uniformly
polynomial-time, length-preserving permutation family $Q_r$, put
\[
 P_r(x,b)=(Q_r(x),b).
\]
Unconstrained truncated collisions are trivial:
$(x,0),(x,1)$. Requiring their common output to equal a supplied label
$y$ makes compatible collision search exactly inversion of $Q_r$ at
target $y$, meaning finding $Q_r^{-1}(y)$.
A recovered inverse gives the pair, and either member of a valid pair
contains that inverse.

The inverse-bit language for a \emph{fixed uniform permutation family}
lies in $\mathsf{UP}\cap\mathsf{coUP}$: the unique preimage certifies
either bit value. For an arbitrary circuit supplied as input, this
statement instead needs the promise that the circuit is a permutation.
Worst-case one-way permutations and the related class separation are
studied by Homan and Thakur \cite{homan2003}.
No superpolynomial inversion lower bound is proved here.
Under $\mathsf P=\mathsf{NP}$, SAT search in the polynomially verifiable
inverse relation gives a polynomial inverter.
If $y$ is constructed as $Q_r(w)$ from an already available witness $w$,
replay already has its preimage.

\section{Explicit constraints and semantic SNSAT must be distinguished}
An explicitly supplied existential Boolean network with $G$ bounded-fan-in
gates and constant-arity constraints has a Tseitin encoding with $O(G)$ variables/clauses and
$O(G\log G)$ bits. A hypothetical SAT algorithm of time $M^a$ solves it
in polynomial time; witness self-reduction uses at most $O(G)$
restricted queries. Circuit depth does not multiply $a$ once this
whole explicit network is encoded.

Ordinary SNSAT instead defines semantic truth bits
\[
 y_i=1\quad\Longleftrightarrow\quad
 \exists z_i\;\phi_i(y_1,\ldots,y_{i-1},z_i).
 \tag{19}
\]
Deciding the final bit is the standard $\mathsf P^{\mathsf{NP}}$-complete
SNSAT problem \cite{markey2003}. False bits must be actual negative SAT answers.
Merely guessing a false bit and one falsifying assignment does not
enforce (19).

If a complete explicit list has total encoded length $N$ and $k$ stages,
each restricted query has length at most $O(N)$. A SAT algorithm with
a nondecreasing bound $T$ evaluates the list in
\[
 \poly(N)+\sum_{i=1}^kT(O(N_i))
 \ \le\ \poly(N)+kT(O(N)),
 \tag{20}
\]
where $N_i$ is its explicitly charged query length.
Polynomial $T$ therefore stays polynomial for $k\le\poly(N)$;
a fixed quasipolynomial logarithmic degree also stays unchanged.
For succinct recursive or alternating variants, one must account for
the resulting queries or expansion. They are not automatically
covered by either (18) or (20).
The notation $\mathrm{SNSAT}^{\mathrm{SNSAT}}$ alone does not specify
which semantic construction is intended.

\section{Verified computations}
All experiments are bounded finite controls, not measurements of $H$.
Every invocation finished below 30 seconds.
\begin{description}
\item[] \path{experiments/review/compatibility_selector_dp_20261111.py}.
SHA-256:
\texttt{5bc4907bd3e865d221433d9819405e5c7a657ce3a2f51c7a2632ce8fc65b39d5}.
One recorded run took about 1.02 seconds.
It exhausts all 256 subsets of the eight proper signed 3-clauses on
three variables, then checks 128 fixed-seed four/five-variable formulas,
nine hostile controls, and (16). Seed: 20261111.
All 394 cases agree with exhaustive SAT: 385 SAT and nine UNSAT.
The projection checks cover 5,372 source assignments and 273,481
local selector patterns. All 385 reconstructed encoded witnesses verify.
\item[] The degree audit fixes a source model of (16), checks all 32
points of the selected five-dimensional cube, and obtains exactly
$\operatorname{AND}_5$, with one ANF monomial of degree five.
For the unsatisfiable conjunction of all eight proper 3-clauses on
three variables, it enumerates all 6,912 Horn models and finds zero
accepting models.
\item[] \path{experiments/review/compatibility_collision_coordinates_20261111.py}.
SHA-256:
\texttt{068118e72ad3d923e47bca3701f4afd6151ef27c3a3970e094789367343d076e}.
The recorded run took about 0.478 seconds.
It explicitly depends on the already verified collision producer,
SHA-256
\texttt{9197bc468d255e7746f6d2f6c8a421c6186bce887584ef12212c8cd26c6a5ff8}.
It checks a two-stage, four-bit Feistel chain with eight proper
compatibility clauses on its six label bits.
All 256 coordinate strings give distinct valid primary chains;
independent forward enumeration checks 4,352 candidate pairs and
finds the identical 256-chain set.
All four adaptive second-stage keys occur.
\item[] The collision source has 16 satisfying label assignments and
64 compatible ordered primary chains. The compiled proper 3-CNF has
201 variables and 805 clauses.
All 1,024 chain-valid primary/padding completions agree with source
compatibility, with 256 accepted completions.
Eight of the 16 oriented first prefixes are dead.
The all-zero-label first prefix is dead, while labels
$100\,|\,100$ lift directly to a full solution.
\end{description}
Every recorded mismatch count is zero.
These checks do not enumerate all auxiliary assignments independently:
functional gate equivalences, previously truth-checked, determine
the completions. The general statements have separate mathematical
proofs and independent temporary-agent review.

Raw artifacts are under
\path{results/raw/review/2026-11-11_COMPATIBILITY_CORE_RETRY/}:
\path{selector_dp.json}, \path{collision_coordinates.json},
and the exact \path{dead_prefix.cnf}.

\section*{CORPUS REVIEWED}
This retry extends the immediately preceding corpus survey and
20:17 report. It directly reopens the material ATC2 manuscript and
the actual executable compiler, rather than inferring the gadget
from a schematic description. Source blobs:
\begin{description}
\item[] ATC2 manuscript:
\texttt{3809296b0a5f6ea85fc03bf1d8441b04e3017585}.
\item[] ATC2 executable:
\texttt{575a6aec2c49e21452f1d6b28d6aca87a8c66f85}.
\item[] Root manuscript snapshot for the parent continuation:
\texttt{ca1c40a11155442a7ecdf74d0e05db3f83abae9a}.
\end{description}
Governance, both root proof trackers, prior transfer/closure reports,
and targeted specialist history were reviewed in the ongoing run.
Markey's original thesis definitions and theorem were opened for
the SNSAT semantics distinction.
This is a targeted read of material original artifacts, not an
all-byte survey or full-checkout test claim.

\section*{DUPLICATION / CONTRADICTION AUDIT}
The selector compiler and general Horn/affine algorithm are prior
components. The degree dichotomy explicitly strengthens the previous
generic warning about acquiring a restricted-degree bound:
on this concrete image the degree-drop question is exactly UNSAT.
The clause-coverage controller sharpens the generic bound for this
specific encoding; it is an elementary algorithm, not a claimed new
complexity-class separation.
The collision coordinate theorem extends the previous inverse control
to arbitrary global compatibility restrictions and makes the actual
remaining constraint explicit.

No contradiction with $\mathsf P=\mathsf{NP}$ is found.
A genuine lower bound for the degree-decision problem, the compatibility
core, or inversion would transfer to a separation, but none follows
from the numerical degree, rank, or number of dead prefixes.

\section*{WHAT CHANGED RELATIVE TO ROOT MAIN}
The exact degree-transfer statement and specialized $2^m$ controller
can refine the parameter ledger. They close a representation-equivalence
question, not the all-program lower-bound arrow.
The canonical root manuscript ownership regions and previous
timestamped reports are unchanged. This is a Codex handoff for the
synthesis owner.
\textbf{TRACKERS REVIEWED; NO MATERIAL CHANGE} to a load-bearing
OPEN status.

\section*{ADOPTED / REJECTED / RETIRED PRIOR COMPONENTS}
\begin{description}
\item[] Adopted: the exact selector compiler, independent source
witness verification, global shared-state computation, and the
fixed-program replay accounting.
\item[] Refined: generic closure degree to the exact gap (10);
seed enumeration to clause-coverage DP on this image;
collision prefixes to global label/orientation coordinates.
\item[] Rejected: high minimum algebraic degree as a runtime lower
bound, dead prefixes as mandatory unrestricted branching,
and freely guessed false SNSAT answers.
\item[] Retired for this architecture: efficiently invertible
collision wrapping with polynomial source-witness lifting as a
fixed-seed growing-$H$ amplifier.
\end{description}

\section*{Hostile controls and barrier audit}
The nine selector controls include empty and contradictory formulas,
short or repeated literals, tautologies and sparse identifiers.
The exact theory charges dense relabelling and only occurring variables.
The connected easy family (11) has arbitrarily high restricted degree.
The explicit dead-prefix family (16) is solved globally.
The collision audit varies all orientation and padding bits.

The algorithms and reductions are model-matched explicit computations.
Oracle or algebraic access cannot invalidate these upper algorithms.
The degree transfer identifies a semantic decision problem; it does
not prove a lower bound against algebraic-oracle machines.
No Natural Proofs barrier is claimed to have been crossed.
No cryptographic hardness hypothesis, ETH assumption, or claimed
ETH refutation is used.

\section*{Status table}
\begin{longtable}{p{0.58\textwidth}p{0.30\textwidth}}
\toprule
Claim & Status\\
\midrule
Exact selector witness projection & PROVED/EXACT; recovered\\
Degree dichotomy (10) and coNP transfer & PROVED/EXACT\\
Clause-coverage algorithm & PROVED/EXACT\\
Actual dead-prefix examples & PROVED/EXACT and COMPUTATIONAL\\
Collision coordinate bijection & PROVED/EXACT\\
Pinned collision/inversion equivalence & PROVED/EXACT\\
Numerical degree/rank implies superpolynomial $H$ & FALSIFIED/RETIRED\\
Growing-$H$ inequality (3) & OPEN\\
P versus NP & OPEN\\
\bottomrule
\end{longtable}

\section*{Single next target}
Use a fully specified global compatibility core and seek a lower bound
on its \emph{uniform computational cost}, with source recovery and
witness maps charged. The exact degree gap supplies one such
decision problem, while the clause-coverage controller supplies an
upper bound every proposed lower bound must respect.
A proof must exclude global compression of states, not merely exhibit
many prefixes or a large algebraic degree.

\bibliographystyle{plain}
\bibliography{research/agents/codex/2026-11-11_2035_COMPATIBILITY_CORE_AND_DEGREE_TRANSFER_refs}
% END_OF_VERIFIED_REPORT::CODEX::COMPATIBILITY-CORE-DEGREE-TRANSFER
\end{document}
